\documentclass[11pt]{amsart}
\usepackage{calc,amssymb,amsmath,ulem,stmaryrd,fullpage}
\usepackage{enumerate}
\usepackage{alltt}
\RequirePackage[dvipsnames,usenames]{color}
\let\mffont=\sf
\normalem
%\input{kmacros3.sty}
\input{xy}
\xyoption{all}
\usepackage{hyperref}
\usepackage{graphicx}
\usepackage[all,cmtip]{xy}
\usepackage{verbatim}
\usepackage[left=0.95in,top=0.95in,right=0.95in,bottom=0.95in]{geometry}
\newcommand{\tauCohomology}{T}
\newcommand{\FCohomology}{S}


%\theoremstyle{theorem}
%\newtheorem*{mainthm*}{Main Theorem}

\newcommand{\note}[1]{\marginpar{\sffamily\tiny #1}}

\def\todo#1{\textcolor{Mahogany}%
{\sffamily\footnotesize\newline{\color{Mahogany}\fbox{\parbox{\textwidth-15pt}{#1}}}\newline}}

\renewcommand{\O}{\mathcal O}
\newcommand{\ie}{{\itshape i.e.} }
\newcommand{\roundup}[1]{\lceil #1 \rceil}
\newcommand{\rounddown}[1]{\lfloor #1 \rfloor}
\renewcommand{\to}[1][]{\xrightarrow{\ #1\ }}
\newcommand{\onto}[1][]{\protect{\xrightarrow{\ #1\ }\hspace{-0.8em}\rightarrow}}
\newcommand{\into}[1][]{\lhook \joinrel \xrightarrow{\ #1\ }}
%\newcommand{DBDual}[1]{{\underline{\omega}_{#1}}}
\newcommand{\DBDual}[1]{{{\uuline \omega}{}^{\mydot}_{#1}}}
\newcommand{\Tt}{{\mathfrak{T}}}
%\newcommand{\bn}{{\mathfrak{n}} }
\newcommand{\sol}[1]{ \vskip 6pt{\color{blue} {\bf Solution:} #1} \vskip 6pt}

\newcommand{\bR}{{\mathbb{R}}}
\newcommand{\va}{{\vec{a}}}
\newcommand{\vb}{{\vec{b}}}
\newcommand{\vzero}{{\vec{0}}}
\newcommand{\vi}{{\vec{i}}}
\newcommand{\vj}{{\vec{j}}}
\newcommand{\vk}{{\vec{k}}}
\newcommand{\vh}{{\vec{h}}}
\newcommand{\vr}{{\vec{r}}}
\newcommand{\vu}{{\vec{u}}}
\newcommand{\vv}{{\vec{v}}}
\newcommand{\vw}{{\vec{w}}}
\newcommand{\vx}{{\vec{x}}}
\newcommand{\vy}{{\vec{y}}}
\newcommand{\vz}{{\vec{z}}}
\newcommand{\vT}{{\vec{T}}}
\newcommand{\vN}{{\vec{N}}}
\newcommand{\vF}{{\vec{F}}}
\newcommand{\vG}{{\vec{G}}}
\newcommand{\bF}{\mathbb{F}}
\newcommand{\bQ}{\mathbb{Q}}
\newcommand{\bZ}{\mathbb{Z}}
\newcommand{\Inn}{\mathrm{Inn}}
\newcommand{\Aut}{\mathrm{Aut}}

\begin{document}
\title{HW \#2 -- Math 6310\\Fall 2019}
\author{Due: Friday, September 13th}
\maketitle

\begin{enumerate}
    \item Recall that the \emph{index} of a subgroup $H \leq G$ is $[G : H]$, the number of cosets (left = right) of $H$.  Show that every subgroup of index $2$ is normal.  Conclude that $A_n$ is normal in $S_n$.
    \newpage
    \item Suppose that $H$ and $H'$ are simple groups\footnote{Groups with no proper nontrivial normal subgroups.} and set $G = H \times H'$ a group with the product (componentwise) group structure.  Show that every normal subgroup of $G$ that is not $\{1\}$ or $G$,is isomorphic to $H$ or $H'$.  \vskip 3pt
        \emph{Hint:} Show that the intersection of normal subgroups is normal.  Note, that not every subgroup will be of the form $H \times \{1\}$ and $\{1\} \times H'$ (in particular, think of the case where $H = \bZ/p\bZ = H'$ under addition).
    \newpage
    \item Compute the cosets of $H = \langle (12) \rangle$ in $S_3$.  Show explicitly with an example that addition of cosets via the formula $(a H) (b H) = (ab) H$ is \emph{not} well defined.
    \newpage
    \item   Suppose $S \subseteq G$ is a subset of a group such that $g S g^{-1} \subseteq S$ for all $g \in G$.  Show that $\langle S \rangle$ is normal.  Next let $T \subseteq G$ be another subset and form the set $V = \bigcup_{g \in G} g^{-1} T g$.  Show that $\langle V \rangle$ is the unique smallest normal subgroup of $G$ containing $T$.
    \newpage
    \item Determine $\Aut S_3$.  Here $\Aut(G)$ is the \emph{group} of bijective group homomorphisms $\phi : G \to G$ (isomorphisms) under composition.
    \newpage
    \item For any $a \in G$ consider the map $\phi_a : G \to G$ defined by $\phi_a(x) = axa^{-1}$.
    \begin{enumerate}
        \item Show that $\phi_a$ is an automorphism.  (It is called an \emph{inner automorphism}).
        \item  Show that $a \mapsto \phi_a$ gives us a homomorphism $G \to \Aut(G)$ with kernel equal to $Z(G)$  the center\footnote{Those elements of $G$ that commute with every other element of $G$} of $G$.
        \item  Let $\Inn(G)$ be the image of the map in (ii).  Show that $\Inn(G) \cong G/Z(G)$.
        \item  Finally show that $\Inn(G)$ is a normal subgroup of $\Aut(G)$.
    \end{enumerate}
    \newpage
    \item {\bf Extra credit: }
    Consider the free group $F$ on the set $\{x_2, \dots, x_{n} \}$.  Consider the normal subgroup $K$ generated by the following elements $x_i^2, (x_i x_j)^3, (x_i x_j x_i x_k)^2$ (for $i,j,k$ distinct).  Show that $F/K$ is isomorphic to $S_n$.
\end{enumerate}

\end{document}

